\subsection{Kufijtë analitikë të modelit Ising}
Në temperaturë zero dhe $H=0$, energjia minimizohet kur të gjitha spinet janë të njëjta. Në rrjetin katror çdo spin ka katër fqinjë, por çdo lidhje numërohet një herë; ka $2N$ lidhje. Prandaj energjia bazë është
\begin{equation}
 E_0=-2JN,qquad e_0=-2J.
\end{equation}
Magnetizimi për spin ka modul një. Këto janë teste të drejtpërdrejta për kodin në temperaturë të ulët.

Në temperaturë të pafundme, konfigurimet janë të barasvlershme dhe spinet janë të pavarura me mesatare zero. Atëherë $\langle E\rangle\to0$, $\langle M\rangle=0$ dhe
\begin{equation}
 \langle M^2\rangle=N.
\end{equation}
Prandaj $\langle M^2\rangle/N^2=1/N$. Një simulim me temperaturë shumë të lartë duhet t'i afrohet këtyre kufijve. Për rrjete shumë të vogla, të gjitha $2^N$ konfigurimet mund të enumerohen dhe të japin vlera të sakta të funksionit të ndarjes; ky është testi më i fortë i algoritmit dhe i observablave.

\subsubsection{Pranimi dhe bilanci energjetik lokal}
Në $H=0$, shuma e katër fqinjëve merr vlerat $-4,-2,0,2,4$. Për rrjedhojë, ndryshimi i energjisë kufizohet në bashkësinë
\begin{equation}
 \Delta E\in\{-8J,-4J,0,4J,8J\}.
\end{equation}
Faktorët $e^{-4\beta J}$ dhe $e^{-8\beta J}$ mund të parallogariten. Ky optimizim shmang eksponencialen në çdo provë dhe ruan saktësisht algoritmin.

Një matje e energjisë duhet të numërojë vetëm lidhjet djathtas dhe lart, ose të ndajë me dy shumën mbi të katër fqinjët. Gabimi i numërimit të dyfishtë ndryshon energjinë dhe temperaturën efektive. Testi lokal zgjedh një spin, llogarit energjinë e plotë para dhe pas përmbysjes dhe kontrollon barazinë me formulën e $\Delta E$ për shumë konfigurime të rastit.

\subsubsection{Kapaciteti termik, susceptibiliteti dhe cumulanti Binder}
Derivimi i pritjes kanonike jep
\begin{equation}
 \frac{\partial\langle E\rangle}{\partial\beta}
 =-\left(\langle E^2\rangle-\langle E\rangle^2\right).
\end{equation}
Me $\beta=1/(k_BT)$ merret formula e fluktuacioneve për $C_V$. Në mënyrë analoge, derivimi ndaj fushës jep susceptibilitetin nga fluktuacionet e magnetizimit. Këto marrëdhënie lidhin përgjigjen makroskopike me luhatjet mikroskopike.

Cumulanti Binder për magnetizimin është
\begin{equation}
 U_L=1-\frac{\langle M^4\rangle}{3\langle M^2\rangle^2}.
\end{equation}
Në fazën e çrregullt Gauss, $U_L\to0$; në një shpërndarje me dy maja të ngushta pranë $\pm M_0$, $U_L\to2/3$. Kurbat për madhësi të ndryshme rrjeti priten pranë temperaturës kritike, duke dhënë një vlerësim më robust se maksimumi i një observable të vetëm.

\subsubsection{Shkallëzimi me madhësinë e fundme}
Pranë pikës kritike, gjatësia e korrelacionit sillet si $\xi\sim|t|^{-\nu}$, ku $t=(T-T_c)/T_c$. Në rrjet të fundëm, rritja ndalet kur $\xi\sim L$. Kjo çon në forma shkallëzimi, për shembull
\begin{equation}
 m(T,L)=L^{-\beta_m/\nu}\mathcal M(tL^{1/\nu}),
\end{equation}
ku $\beta_m$ është eksponent kritik, jo inversi i temperaturës. Në $T_c$, grafiku log--log i $m$ ndaj $L$ ka pjerrësi $-\beta_m/\nu$.

Maksimumi i susceptibilitetit rritet afërsisht si $L^{\gamma/\nu}$ dhe temperatura e pseudokritike zhvendoset si $L^{-1/\nu}$. Një studim serioz përdor disa madhësi, interpolon observablat me pacaktueshmëri dhe kontrollon shembjen e të dhënave. Një kurbë e vetme magnetizimi nuk mjafton për të përcaktuar tranzicionin.

\subsubsection{Jackknife për madhësi jolineare}
Le të ndahen matjet në $B$ blloqe. Për çdo $b$ llogaritet statistiku $\theta_{(b)}$ duke hequr bllokun $b$. Mesatarja jackknife është $\overline\theta_J=B^{-1}\sum_b\theta_{(b)}$ dhe varianca vlerësohet me
\begin{equation}
 \widehat{\Var}_J(\theta)=\frac{B-1}{B}
 \sum_{b=1}^{B}(\theta_{(b)}-\overline\theta_J)^2.
\end{equation}
Kjo metodë përhap kovariancën ndërmjet $\langle E\rangle$ dhe $\langle E^2\rangle$ në $C_V$, ose ndërmjet momenteve në cumulantin Binder. Blloqet duhet të jenë më të gjata se autokorrelacioni; përndryshe pseudo-mostrat nuk janë afërsisht të pavarura.

\subsubsection{Krahasimi i algoritmeve}
Kostoja nuk matet vetëm me sweeps. Krahasimi përdor kohën e procesorit për një mostër efektivisht të pavarur,
\begin{equation}
 C_{\mathrm{eff}}\propto C_{\mathrm{sweep}}\,2\tau_{\mathrm{int}}.
\end{equation}
Metropolis-i lokal ka sweep të lirë, por $\tau$ të madh pranë kritikës. Algoritmi Wolff ndërton një grup spinesh të lidhura dhe e përmbys; kostoja për hap ndryshon, por autokorrelacioni kritik zvogëlohet fort. Krahasimi duhet bërë për të njëjtin observable dhe të njëjtën madhësi rrjeti.

Koha Monte Carlo e këtyre algoritmeve nuk është dinamika reale e magnetit, përveç rasteve kur rregullat e përditësimit janë zgjedhur si model kinetik. Për observablat e ekuilibrit mjafton shpërndarja stacionare; për koha relaksimi fizik nevojitet model dinamik me interpretim të veçantë.
